\section{Related Work}
\label{sec:related}

\paragraph{Deep research systems.}
OpenAI Deep Research \citep{openai2025deepresearch} is an end-to-end agent that plans searches, browses, and writes a cited report. It was trained with reinforcement learning on browsing tasks of the kind collected in BrowseComp \citep{wei2025browsecomp}, where answers are short facts that are hard to find, so its reward favors retrieval and citation accuracy over the structure of an argument. Its reasoning lives in one context window with no persistent dependency structure between sub-questions. STORM \citep{shao2024storm} writes Wikipedia-like articles: it discovers perspectives on a topic, simulates conversations in which writers with those perspectives question a search-grounded expert, curates the answers into an outline, and writes the article. The questions emerge from the conversations as the research proceeds, so the process is bottom-up, and the output is deliberately balanced. Both systems were designed for summary, and the evaluation below treats them as the strongest available summary-style baselines, not as argumentation systems.

\paragraph{Decomposition and structured reasoning.}
Breaking a problem into sub-problems before solving it is an established prompting strategy: least-to-most prompting \citep{zhou2023leasttomost}, decomposed prompting \citep{khot2023decomposed}, and self-ask \citep{press2023selfask} all have a model generate sub-questions whose answers feed the original question. Tree of Thoughts \citep{yao2023tot} searches over intermediate reasoning steps, ReAct \citep{yao2023react} interleaves reasoning with tool calls, and Graph of Thoughts \citep{besta2024got} generalizes the tree to a graph whose thoughts can merge. These methods target problems with checkable answers, mostly arithmetic, puzzles, or short-form QA. They do not synthesize prose from retrieved evidence, and they do not refine the reasoning structure when a composed answer turns out to be incomplete. \sys{} borrows the decomposition idea and adds composition instructions, evidence-grounded leaves, and gap-driven refinement of the graph.

\paragraph{Multi-agent frameworks.}
AutoGen \citep{wu2023autogen}, CAMEL \citep{li2023camel}, and natural-language societies of mind \citep{zhuge2023mindstorms} coordinate several language-model agents through conversation, with roles such as planner, researcher, and critic. Their communication is flat: agents exchange messages, and the structure of the task is not represented. \sys{} fixes the structure explicitly in the DAG, so that the output of one sub-task constrains the composition of its parent.

\paragraph{Argument mining and generation.}
Argument mining identifies claims, premises, and the relations between them in existing text \citep{lawrence2020argmining}. Argument generation produces new arguments: \citet{hua2018argument} generate counter-arguments with a sequence-to-sequence model conditioned on retrieved evidence, and Project Debater \citep{slonim2021debater} assembles a debate speech from a mined corpus of claims and evidence. These systems draw on fixed corpora and produce a single pass of text. \sys{} researches the web live and builds the argument structure first, then fills it.

\paragraph{Citation and attribution.}
Systems that answer with citations can be checked for whether each citation supports its sentence. \citet{rashkin2023ais} define attribution to identified sources as the standard. \citet{liu2023verifiability} applied human judgments to four generative search engines and found that only about half of generated sentences were fully supported by their citations and about three quarters of citations supported their sentence. \citet{gao2023alce} built a benchmark that scores citation precision and recall automatically with an entailment model. These studies evaluate single-step generation from retrieved passages. \Cref{sec:citations} looks at a multi-step setting, where citations must survive several rounds of composition, and uses a lexical test that needs no model.

\paragraph{Implementation.}
The pipeline is written with DSPy \citep{khattab2024dspy}, which separates each language-model call into a signature (inputs, outputs, and an instruction) and lets the calls be composed as a program. The signatures are listed in \Cref{app:prompts}.
